% Fuentes: tab_s07_transitorios_red.csv
\begin{table}[htbp]
\centering
\caption{Intervalos de los seis estados en la prueba numérica de la red térmica.}
\label{tab:res_red_estados}
\small
\setlength{\tabcolsep}{5pt}
\begin{tabular}{lrrrr}
\toprule
Estado & Unidad & Mínimo & Máximo & Media \\
\midrule
$T_c$ & K & $298{,}150000$ & $300{,}480314$ & $299{,}577756$ \\
$T_b$ & K & $298{,}150000$ & $300{,}171307$ & $299{,}213025$ \\
$T_p$ & K & $298{,}076397$ & $299{,}005635$ & $298{,}416732$ \\
$T_a$ & K & $298{,}150000$ & $299{,}331106$ & $298{,}718578$ \\
$W$ & kg vapor/kg aire seco & $0{,}007861$ & $0{,}008604$ & $0{,}008369$ \\
$M$ & kg agua/kg materia seca & $0{,}094780$ & $0{,}100000$ & $0{,}097328$ \\
\bottomrule
\end{tabular}
\par\vspace{3pt}
\begin{minipage}{0.98\textwidth}\footnotesize
Prueba de software de $1200$ s con salidas programadas cada $10$ s. Los parámetros y las condiciones corresponden al caso numérico, sin identificación experimental de la planta.
\end{minipage}
\end{table}
